
\documentclass{article}
\usepackage{colortbl}
\usepackage{makecell}
\usepackage{multirow}
\usepackage{supertabular}

\begin{document}

\newcounter{utterance}

\centering \large Interaction Transcript for game `cladder', experiment `full\_v1.5\_default', episode 8830 with Bentel rockers\_\_qwen3.5{-}4b\_\_Bentel\_iter.
\vspace{24pt}

{ \footnotesize  \setcounter{utterance}{1}
\setlength{\tabcolsep}{0pt}
\begin{supertabular}{c@{$\;$}|p{.15\linewidth}@{}p{.15\linewidth}p{.15\linewidth}p{.15\linewidth}p{.15\linewidth}p{.15\linewidth}}
    \# & \multicolumn{2}{c}{Player} && \multicolumn{2}{c}{Game Master} \\
    \hline

    \theutterance \stepcounter{utterance}  
    & & & \multicolumn{4}{p{0.6\linewidth}}{
        \cellcolor[rgb]{0.9,0.9,0.9}{
            \makecell[{{p{\linewidth}}}]{
                \texttt{\tiny{[P1$\langle$GM]}}
                \texttt{You are an expert in causal inference. The following question is not a typical commonsense query, but rather a meticulously designed question created by a professor specializing in causal inference, intended to assess the students' mastery of the course content.} \\
\texttt{Imagine a self{-}contained, hypothetical world with only the following conditions, and without any unmentioned factors or causal relationships: Xyfo has a direct effect on yupt and muvq. Zupj has a direct effect on yupt. Yupt has a direct effect on muvq. Xyfo is unobserved. Method 1: We look at how yupt correlates with muvq case by case according to xyfo. Method 2: We look directly at how yupt correlates with muvq in general. To understand how yupt affects muvq, is it more correct to use the Method 1 than Method 2?} \\
\texttt{Start your answer with 'Yes' or 'No', followed by additional reasoning or evidence to support your explanation.} \\
            }
        }
    }
     \\ \\

    \theutterance \stepcounter{utterance}  
    & \multicolumn{4}{p{0.6\linewidth}}{
        \cellcolor[rgb]{0.561,0.737,0.561}{
            \makecell[{{p{\linewidth}}}]{
                \texttt{\tiny{[P1$\rangle$GM]}}
                \texttt{Yes} \\
\\ 
\texttt{**Reasoning:**} \\
\\ 
\texttt{In this hypothetical world, we are dealing with a classic **confounding** scenario. The causal structure is as follows:} \\
\texttt{1.  \$X\_\{yfo\} \textbackslash{}rightarrow Y\_\{upt\}\$} \\
\texttt{2.  \$X\_\{yfo\} \textbackslash{}rightarrow M\_\{uvq\}\$} \\
\texttt{3.  \$Z\_\{upj\} \textbackslash{}rightarrow Y\_\{upt\}\$} \\
\texttt{4.  \$Y\_\{upt\} \textbackslash{}rightarrow M\_\{uvq\}\$} \\
\\ 
\texttt{Here, \$X\_\{yfo\}\$ acts as a **confounder** because it has direct causal effects on both the exposure (\$Y\_\{upt\}\$) and the outcome (\$M\_\{uvq\}\$). Additionally, there is a direct causal pathway from \$Y\_\{upt\}\$ to \$M\_\{uvq\}\$.} \\
\\ 
\texttt{*   **Analysis of Method 2 (General Correlation):** If we simply look at the correlation between \$Y\_\{upt\}\$ and \$M\_\{uvq\}\$ without conditioning on \$X\_\{yfo\}\$, we observe a mix of two effects:} \\
\texttt{    1.  The direct causal effect of \$Y\_\{upt\}\$ on \$M\_\{uvq\}\$.} \\
\texttt{    2.  The spurious correlation induced by the common cause \$X\_\{yfo\}\$.} \\
\texttt{    Because \$X\_\{yfo\}\$ is unobserved, we cannot statistically adjust for it (e.g., via regression or matching) to isolate the direct effect. Therefore, Method 2 yields a biased estimate of the causal effect of \$Y\_\{upt\}\$ on \$M\_\{uvq\}\$.} \\
\\ 
\texttt{*   **Analysis of Method 1 (Stratification/Conditioning):** By examining the correlation between \$Y\_\{upt\}\$ and \$M\_\{uvq\}\$ *case by case* according to the levels of \$X\_\{yfo\}\$, we effectively condition on the confounder.} \\
\texttt{    *   Within strata where \$X\_\{yfo\}\$ is constant, the variation in \$Y\_\{upt\}\$ is not driven by \$X\_\{yfo\}\$.} \\
\texttt{    *   Consequently, the observed association within these strata reflects only the direct causal pathway (\$Y\_\{upt\} \textbackslash{}rightarrow M\_\{uvq\}\$), as the confounding bias introduced by \$X\_\{yfo\}\$ is removed.} \\
\\ 
\texttt{Since \$X\_\{yfo\}\$ is the only unobserved factor mentioned and it is a confounder, Method 1 (stratification) is the correct approach to isolate the true causal effect of \$Y\_\{upt\}\$ on \$M\_\{uvq\}\$ in this specific closed{-}world scenario. This process is known as **backdoor adjustment** (specifically, blocking the backdoor path \$Y\_\{upt\} \textbackslash{}leftarrow X\_\{yfo\} \textbackslash{}rightarrow M\_\{uvq\}\$).} \\
            }
        }
    }
    & & \\ \\

    \theutterance \stepcounter{utterance}  
    & & & \multicolumn{2}{p{0.3\linewidth}}{
        \cellcolor[rgb]{0.9,0.9,0.9}{
            \makecell[{{p{\linewidth}}}]{
                \texttt{\tiny{[GM$|$GM]}}
                \texttt{Yes} \\
            }
        }
    }
    & & \\ \\

    \theutterance \stepcounter{utterance}  
    & & & \multicolumn{2}{p{0.3\linewidth}}{
        \cellcolor[rgb]{0.9,0.9,0.9}{
            \makecell[{{p{\linewidth}}}]{
                \texttt{\tiny{[GM$|$GM]}}
                \texttt{yes} \\
            }
        }
    }
    & & \\ \\

    \theutterance \stepcounter{utterance}  
    & & & \multicolumn{2}{p{0.3\linewidth}}{
        \cellcolor[rgb]{0.9,0.9,0.9}{
            \makecell[{{p{\linewidth}}}]{
                \texttt{\tiny{[GM$|$GM]}}
                \texttt{game\_result = WIN} \\
            }
        }
    }
    & & \\ \\

\end{supertabular}
}

\end{document}
